% ============================================================
% ICÔNE OMARCHY — SIDEBAR
% ============================================================
% Le dessin lui-même est dans config/icon.tex (partagé avec la variante une
% colonne, où il sert de filigrane centré sur chaque page).
%
% Placement : l'icône est simplement dans le flux, encadrée par les deux \vfill
% de sections/sidebar.tex, ce qui la centre dans l'espace libre sous le contenu
% de la sidebar. Volontairement PAS d'overlay « remember picture » : celui-ci
% demanderait deux passes de compilation, alors que le README documente un
% simple « pdflatex main.tex » — au premier passage l'icône serait mal placée.
%
% \omarchyiconsize : côté de l'icône (carrée). Borné par \linewidth, la largeur
% utile de la sidebar.
\newlength{\omarchyiconsize}
\setlength{\omarchyiconsize}{3.4cm}

% Teinte très diluée : filigrane à peine perceptible. On mélange la couleur au
% fond de la sidebar plutôt que d'utiliser « opacity », ce qui donne exactement
% le même rendu ici (le fond est un aplat uni) mais reste une couleur opaque
% ordinaire — pas de canal de transparence, donc aucun risque à l'impression ni
% dans les lecteurs PDF qui aplatissent mal les calques.
% \omarchyiconmix : part de maincolor, en pourcentage (0 = invisible).
\newcommand{\omarchyiconmix}{12}

{\centering
\omarchyicon{\omarchyiconsize}{maincolor!\omarchyiconmix!sidebarcolor}
\par}
